\section{Challenges and Countermeasures}
\label{sec:challenges}

Building a clean, reproducible Ethereum fraud detector from public data
sources requires overcoming a set of recurring engineering and
modelling challenges.  This section maps each challenge to the
architectural or procedural choice that directly addresses it,
connecting the practical decisions to the limitations identified in
prior work (Section~\ref{sec:background}).

Blockchain's tamper-proof and auditable properties have been leveraged for
secure provenance management in distributed
systems~\cite{poudel2025sphere}.  The same auditability principles guide
our pipeline design: every address transformation, join, and label
assignment is deterministic and reproducible from the raw source files.

\subsection{Challenge 1: Address Format Inconsistency}
Ethereum addresses appear in lowercase in BigQuery exports, in EIP-55
mixed-case checksums in Forta exports, and in either form in Etherscan
API responses.  Silent join failures caused by this inconsistency
corrupted our earliest dataset attempts, producing false negatives where
known fraud addresses failed to match their labels.

\textbf{Countermeasure.}  Every address is lowercased on read at every
collector entry point, implemented once in
\texttt{src/preprocessing/addresses.py} and reused by every pipeline
module.  This eliminates case-induced mismatches and whitespace artefacts
that occasionally appear in CSV exports, achieving 100\% join consistency
across all three data sources.

\subsection{Challenge 2: Sparse Verified Source Code}
Fewer than 0.35\% of the \num{85680} addresses in our transaction graph
carry verified Solidity source.  A pipeline that requires source for every
node would discard almost the entire graph and strip out the structural
signals that the GNN needs.

\textbf{Countermeasure.}  Nodes enter the graph with an empty
\texttt{contract\_text} field when no source is available.  The graph
encoder processes \emph{all} nodes during neighbourhood aggregation,
preserving relational information.  The fusion classifier is trained only
on the subset of nodes that carry both source code and an explicit label
(304 contracts), while the full graph neighbourhood of 85,680 nodes
informs the GraphSAGE embeddings.

\subsection{Challenge 3: Etherscan API Rate Limits}
The Etherscan free tier permits approximately five requests per second.
Exceeding this limit produces HTTP 429 errors and corrupts the source-code
dataset.

\textbf{Countermeasure.}  A 250\,ms sleep is enforced between successive
API calls, and every request is wrapped in a try/except block that logs
and skips on failure without interrupting the pipeline.  Source retrieval
for 660 seed addresses completes in approximately 30 minutes.

\subsection{Challenge 4: Class Imbalance}
The labelled subset contains 95 fraud contracts out of 304 (31.3\%).
Training under unweighted cross-entropy drives the model toward the
majority class, yielding near-zero fraud recall.

\textbf{Countermeasure.}  Weighted cross-entropy (Equation~\ref{eq:loss})
assigns penalty weight $w_1 = n_\text{neg}/n_\text{pos} \approx 2.37$ to
fraud misclassifications.  This directly addresses the class-imbalance
limitation identified by Chen et al.~\cite{chen2024dahgnn} and ensures
that the minority class contributes proportionally to the training signal.

\subsection{Challenge 5: Embedding Dimensionality Mismatch}
The graph encoder produces 64-dimensional embeddings while CodeBERT
produces 768-dimensional embeddings.  Combining representations of such
different sizes risks one modality dominating the other.

\textbf{Countermeasure.}  We concatenate the two embeddings to form an
832-dimensional vector and project through a learned linear layer to a
128-dimensional bottleneck (Equation~\ref{eq:fusion}).  The bottleneck
forces the optimiser to discover the relative weighting of the two
modalities rather than fixing it by hand, which is how the architecture
addresses the insufficient cross-modal integration identified as a
limitation in prior hybrid methods~\cite{sun2025tlmg4eth}.

\subsection{Challenge 6: Compute Budget for Large Language Models}
Training the fusion model end-to-end with CodeBERT unfrozen would not
fit within the memory of a single consumer GPU.

\textbf{Countermeasure.}  All 125M CodeBERT parameters are frozen, and
only the 124K parameters of the GraphSAGE encoder and fusion head are
trained.  This reduces memory by two orders of magnitude while
retaining the full pre-trained code semantics, stabilises training by
eliminating gradient flow through a model that was not designed for
fine-tuning under weighted loss, and enables each training epoch to
complete in $\approx$34 seconds on CPU.

\subsection{Challenge 7: BigQuery Result Volume}
A na\"ive query for all transactions involving 660 seed addresses
returns $>$665 million rows (approximately 130\,GB), which exceeds the
1\,GB single-file export limit and is impractical to download and process
locally.

\textbf{Countermeasure.}  We apply a balanced sampling strategy: keep the
2,000 most-recent transactions per fraud address and the 300 most-recent
per benign address.  This yields a manageable 239,742-row dataset
(approximately 41\,MB) that preserves fraud signal while preventing
high-volume benign contracts (WETH, USDT) from overwhelming the graph
structure.
